\documentclass{birkjour}
\usepackage{graphicx}
\usepackage{amsmath}
\usepackage{amssymb}
\usepackage{url}
\usepackage[dvipsnames]{color}

% extended math fonts
\usepackage{yhmath}

\usepackage{srctex} % MikTex,Inverse Search in tex files from Yap screen


\newcommand{\e}[1]{\mathbf{e}_{#1}} % Clifford basis vectors in Cl(3,1)
\newcommand{\Ie}[1]{I\negthinspace\mathbf{e}_{#1}}%bivector in Cl(3,1)
\newcommand{\inv}[1]{\overline{#1}} %space inversion
%\newcommand{\ba}{\mathbf{a}} %bold a
\newcommand{\bk}{\mathbf{k}}
\newcommand{\bx}{\mathbf{x}}
\newcommand{\bbf}{\mathbf{f}} %bold f


\def\m#1{\mathsf{#1}} % general multivector

\def\sR{\mathsf{R}} %sans serif Rotation
\def\sH{\mathsf{H}} %sans serif Hamiltonian
\def\sF{\mathsf{F}} %sans serif F
\def\sG{\mathsf{G}} %sans serif G


\newcommand{\cl}[2]{\ensuremath{\mathit{Cl}_{#1,#2}}}

% Real, complex and quaternion fields
\newcommand{\bbR}{\ensuremath{\mathbb{R}}}
\newcommand{\bbC}{\ensuremath{\mathbb{C}}}
\newcommand{\bbH}{\ensuremath{\mathbb{H}}}

%involutions
\newcommand{\reverse}[1]{\widetilde{#1}}
\newcommand{\gradeinverse}[1]{\wideparen{#1}}
\newcommand{\cliffordconjugate}[1]{\widetilde{\wideparen{#1}}}

% Quaternion base vectors
\newcommand{\ii}{\mathrm{i}}
\newcommand{\jj}{\mathrm{j}}
\newcommand{\kk}{\mathrm{k}}
%\newcommand{\ii}{\mathrm{i}} %imaginary unit
\newcommand{\ee}{\mathrm{e}} %exponent
\newcommand{\dd}{\mathrm{d}} %differential



% General multivectors
\def\A{\mathsf{A}}
\def\B{\mathsf{B}}
\def\C{\mathsf{C}}
\def\D{\mathsf{D}}
\def\X{\mathsf{X}}
\def\m#1{\mathsf{#1}}%general multivector Input: $\m{a}$

% base vectors and related
\def\e#1{\mathbf{e}_{#1}} %single numbered Clifford basis vectors. Input: $\e3$ or $\e{12}$
\def\g#1{\gamma_{#1}} %single numbered relativistic basis vectors. Input: $\g3$ or $\g{12}$

% vectors = bold lower case
\newcommand{\ba}{\ensuremath{\mathbf{a}}}
\newcommand{\bb}{\ensuremath{\mathbf{b}}}
\newcommand{\bc}{\ensuremath{\mathbf{c}}}
\newcommand{\bd}{\ensuremath{\mathbf{d}}}
\newcommand{\be}{\ensuremath{\mathbf{e}}}
\newcommand{\bv}{\ensuremath{\mathbf{v}}}
\newcommand{\bV}{\ensuremath{\mathbf{V}}}


% bivectors = calligraphic
\newcommand{\cA}{\ensuremath{\mathcal{A}}}
\newcommand{\cB}{\ensuremath{\mathcal{B}}}
\newcommand{\cC}{\ensuremath{\mathcal{C}}}

\newcommand{\magnitude}[1]{\lvert #1\rvert}
\newcommand{\norm}[1]{\lVert #1\rVert}

%text coloring
\newcommand{\blue}[1]{\textcolor{blue}{#1}}
\newcommand{\red}[1]{\textcolor{red}{#1}} % text coloring

\begin{document}

%Insert here the title, affiliations and abstract:

\title[Square root of a multivector of Clifford algebras in~3D]
 {Square root of a multivector of Clifford algebras in~3D: A game with signs}

%----------Author 1
\author{A.~Acus}
\address{Institute of Theoretical Physics and Astronomy,\br
Vilnius University,\br Saul{\.e}tekio 3, LT-10257 Vilnius,
Lithuania} \email{arturas.acus@tfai.vu.lt}

%\thanks{This work was completed with the support of our \TeX-pert.}
%----------Author 2
\author{A.~Dargys}

\address{%
Center for Physical Sciences and Technology,\br Semiconductor
Physics Institute,\br Saul{\.e}tekio 3, LT-10257 Vilnius,
Lithuania}

\email{adolfas.dargys@ftmc.lt}
%----------classification, keywords, date
%----------classification, keywords, date
\subjclass{Primary 15A18; Secondary 15A66}

\keywords{{C}lifford algebra, geometric algebra, square root of multivector, computer-aided theory}

\date{June, 2019}
%----------additions
%\dedicatory{To my boss}


\begin{abstract}
Explicit formulas of the square root of a multivector of real
Clifford algebras \cl{3}{0}, \cl{2}{1}, \cl{1}{2} and \cl{0}{3}
are obtained in radicals. We show that there are four isolated
roots in a case of the most general (generic) multivector (MV) in
the algebras \cl{3}{0},  \cl{1}{2} and \cl{0}{3}. The algebra
\cl{2}{1} makes up an exception and the MV here can have up to 16
isolated roots. In special cases, when the MV is shorter or
restrictions (conditions) between  MV coefficients exist, a
continuum of roots (infinity of roots) on parameter hypersurfaces
can appear in addition to the isolated roots. A number of examples
are presented to illustrate properties of various roots that may
appear in $n=3$ Clifford algebras. Noncommutative Riccati
quadratic equation is solved as an illustration of the application
of the described root algorithm.
\end{abstract}

\maketitle



%\tableofcontents

\section{Introduction}\label{sec:intro}
The square root has a long history. Solution by radicals of the
cubic equation was first published  in 1545 by G.~Cardano.
Simultaneously, a concept of square root of a negative number  has
been developed~\cite{Grant2015}. In 1872 A.~Cayley was the first
to carry over the notion of square root to
matrices~\cite{Cayley1872}. In the recent book by  N.~J.
Higham~\cite{Higham08}, where an extensive literature is presented
on nonlinear functions of matrices, two sections are devoted to
the analysis of square root. The existence of the square root of
matrix is related to it's positive eigenvalues. In the context of
Clifford algebra (CA) the main attempts up till now were
concentrated on the square roots of
quaternions~\cite{Niven1942,Opfer2017}, or their derivatives such
as coquaternions (also called split quaternions), or
nectarines~\cite{Falcao2018,Opfer2017,Ozdemir2009}. The square
root of biquaternion (complex quaternion) was considered
in~\cite{Sangwine2006}. The quaternions and related objects are
isomorphic to one of $n=2$ algebras \cl{0}{2}, \cl{1}{1},
\cl{2}{0} and, therefore, the quaternionic square root analysis
can be easily rewritten in  terms of CA. In this paper we shall be
interested in higher, namely, $n=3$ CAs where the main object is
the 8-component MV.


For the algebras of vector space dimension $n=3$, to say nothing
of higher dimensions, the understanding and investigation of
square root is still in infancy. Some formulas are not complete or
even erroneous for the case of a general MV~\cite{Chapell2015}.
The most akin to the present paper are the investigation of
conditions for existence of square root of
$-1$~\cite{Hitzer2011,Hitzer2013,Sangwine2006}. The existence of
such roots allows to extend the notion of Fourier transform to
multivectors, where they are used in formulating Clifford-Fourier
transform and CA based wavelet theories~\cite{Hitzer13B}.

Our preliminary investigation~\cite{Dargys2019} on this subject
were concerned with square roots of individual MV grades such as
scalar, vector, bivector, pseudoscalar, or their simple
combinations. For this purpose we have applied the Gr{\"o}bner
basis algorithm to analyze the  system of nonlinear polynomial
equations that ensue from the MV equation $\A^2=\B$, where $\A$
and $\B$ are the MVs. The Gr{\"o}bner basis is accessible in
symbolic mathematical packages such as \textit{Mathematica} and
\textit{Maple}. Specifically, the \textit{Mathematica} commands
such as \textbf{Reduce[~]}, \textbf{Solve[~]},
\textbf{Eliminate[~]} and others also employ the Gr{\"o}bner basis
to solve nonlinear problems. With the help of them we were able to
find the characteristic properties of roots for $n=3$ case,
namely, that the MV may have no roots, a single or multiple
isolated roots, or even an infinite number of roots in parameter
spaces of 4D or smaller dimensions.

In this paper we elaborate our investigation~\cite{Dargys2019} of
the problem of the square root of general MV in CA for $n=3$ case.
For this purpose a symbolic experimental
package~\cite{AcusDargys2017} based on {\it Mathematica} system
was used which appeared to be invaluable both for detecting
possible solutions of the nonlinear CA equation $\A^2=\B$ and for
extensive numerical check. In Sec.~\ref{sec:notations}, we
introduce the notation.  The procedure to calculate the square
root of generic MV as well as the special cases that follow is
given in Secs.~\ref{sec:Cl30}-\ref{sec:Cl21} for \cl{3}{0},
\cl{0}{3}, \cl{1}{2}, and \cl{2}{1} algebras, respectively. The
algorithm  is illustrated by examples. Finally,
in~Sec.~\ref{sec:application} the application of a square root of
MV to solve Riccati equation in CA form is demonstrated and
Conclusions are drawn in~Sec.~\ref{sec:conclusions}.


\section{Notations}\label{sec:notations}
A general MV in CA for $n=3$ can be expanded in the orthonormal
basis that consists of $2^{n}=8$ elements listed in inverse degree
lexicographic ordering\footnote{\label{note1}Note an increasing
order of indices in the basis elements, i.e., we employ $\e{13}$
instead of $\e{31}=-\e{13}$. This convention is reflected in
opposite signs of some terms in formulas.}
\begin{equation}\label{basisDim3}
\{1,\e{1},\e{2},\e{3},\e{12},\e{13},\e{23},\e{123}\equiv I\},
\end{equation}
where $\e{i}$ are basis vectors and $\e{ij}$ are the bivectors
(oriented planes). The last term is called the pseudoscalar. The
number of subscripts indicates the grade of basis element. The
scalar is a grade-0 element, the vectors $\e{i}$ are the grade-1
elements etc. In the orthonormalized basis the geometric (Clifford) products of basis vectors satisfy the anticommutation relation,
 \begin{equation}\label{anticom}
 \e{i}\e{j}+\e{j}\e{i}=\pm 2\delta_{ij}.
 \end{equation}
 For \cl{3}{0}
and \cl{0}{3} algebras, correspondingly, the squares of basis
vectors are $\e{i}^2=+1$ and $\e{i}^2=-1$, where $i=1,2,3$. For
mixed signature algebras such as  \cl{2}{1} and \cl{1}{2} we have
$\e{1}^2=\e{2}^2=1$, $\e{3}^2=-1$ and $\e{1}^2=1$,
$\e{2}^2=\e{3}^2=-1$, respectively. The sign of squares of higher
grade elements is determined by squares of vectors and the
property~\eqref{anticom}. For example, in \cl{3}{0} we have
$\e{12}^2=\e{12}\e{12}=-\e{1}\e{2}\e{2}\e{1}=-\e{1}(+1)\e{1}=-\e{1}\e{1}=-1$.
However, in \cl{1}{2} the same computation gives
$\e{12}^2=-\e{1}\e{2}\e{2}\e{1}=-\e{1}(-1)\e{1}=\e{1}\e{1}=+1$.

Any MV $\A$ of real CA of $n=3$ can be expanded in the
basis~\eqref{basisDim3},
\begin{equation}\begin{split}\label{mvAA}
\A=&a_0+a_1\e{1}+a_2\e{2}+a_3\e{3}+a_{12}\e{12}+a_{23}\e{23}+a_{31}\e{31}+a_{123}I\\
\equiv&a_0+\ba+\cA+a_{123}I,
\end{split}\end{equation} where $a_i$, $a_{ij}$ and $a_{123}$ are the real
coefficients, and $\ba=a_1\e{1}+a_2\e{2}+a_3\e{3}$ and
$\cA=a_{12}\e{12}+a_{23}\e{23}+a_{31}\e{31}$ is, respectively, the
vector and the bivector. We will seek for a MV $\A$ the square of
which $\A^2$ will match a given MV $\B$,
\begin{equation}\label{AA}
  \B=\A\A\equiv\A^2=b_0+\mathbf{b}+\cB+b_{123}I.
\end{equation}
In this case we will call the MV $\A$ a square root of $\B$, i.e.,
$\A=\sqrt{\B}$.

In Eq.~\eqref{AA} the MV $\A^2$ has been expanded in the
orthonormal basis where $b_0,\mathbf{b},\cB$ and $I\equiv I_3 $
denote, respectively, a scalar, a vector
($\bb=b_1\e{1}+b_2\e{2}+b_3\e{3}$), a bivector
($\cB=b_{12}\e{12}+b_{23}\e{23}+b_{31}\e{31}$) and a pseudoscalar.
However, the commonly assumed representation~\eqref{mvAA}
of MV $\A$ is not convenient for our problem. Therefore, for
algebras \cl{3}{0}, \cl{0}{3}, \cl{1}{2}, and \cl{2}{1} we
introduce a more symmetric representation of MV,
\begin{equation}
\A=s+\bv+(S+\bV)I,\label{mvA30}
\end{equation}
where now both $s$ and $S$ are the real scalars and both
$\bv=v_1\e{1}+v_2\e{2}+v_3\e{3}$ and
$\bV=V_1\e{1}+V_2\e{2}+V_3\e{3}$ are the vectors with real
coefficients $v_i$ and $V_i$. As we shall see, the MV
representation~\eqref{mvA30} allows to disentangle the coupled
nonlinear equations in a regular manner using similar formulas for
all listed algebras. To select the scalar $s$ in~\eqref{mvA30},
the grade selector
$\langle\m{A}\rangle\equiv\langle\m{A}\rangle_0=s$ is used. The
pseudoscalar part  can be extracted by
$\langle\m{A}\rangle\equiv\langle-\m{A}I\rangle_0=S$, and
similarly for other grades. More about CAs and MV properties can
be found, for example, in books~\cite{Doran03,Lounesto97}.



\section{\cl{3}{0} algebra\label{sec:Cl30}}
The section describes the method of
variable substitution which allows to obtain explicit formulas of
square root of a general MV for all CAs of vector space dimension $n=3$. We
start from Euclidean \cl{3}{0} algebra which presents the most
simple case. Computations with algebras having different signature
will take a similar route and will be presented in
Sections~\ref{sec:Cl03}--\ref{sec:Cl21}.

Our goal is to solve the MV equation $\A^2=\B$, where $\B$ is
known and has the general MV form~\eqref{AA}, or in the
expanded form $\B= b_0+b_1 \e{2}+b_2\e{2}+b_3\e{3}+
b_{12}\e{23}+b_{13}\e{13}+b_{23} \e{23}+b_{123} I$,  and where
$\A$ is represented in the form~\eqref{mvA30}. Expanding $\A^2$
and equating coefficients at all basis elements to $\B$ one
obtains  a system of eight nonlinear equations:
\begin{align}
b_0&=s^2-S^2+\bv^2-\bV^2, &b_{123}&=2(sS+\bv\cdot\bV)\label{MVM1Cl30}, \\
b_1&=2(s v_1-S V_1), &b_{23}&=2(s V_1+S v_1)\label{MVM2Cl30}, \\
b_2&=2(s v_2-S V_2), &b_{13}&=2(s V_2+S v_2)\label{MVM3Cl30}, \\
b_3&=2(s v_3-S V_3), &b_{12}&=2(s V_3+S v_3)\label{MVM4Cl30}.
\end{align}
In the progress of analysis it will be shown that for all CAs in
the consideration the square root algorithm splits into two cases:
the generic case when $s^2 \ne S^2$ and the special case when
$s^2=S^2$. The special case in turn is a union of three subcases:
1) $s = S\neq 0$,  2) $s = -S\neq 0$ and 3) $s=S=0$. We will start
from the generic case.

\subsection{The generic case, $s^2 \neq S^2$}
The system of six Eqs. \eqref{MVM2Cl30}-\eqref{MVM4Cl30} is linear
in $v_i$ and $V_i$. It has  very simple solution which is a key to
analysis that follows\footnote{Note, that the symmetry of the
expression with respect to the coefficients  ($v_2,V_2$) and the
($v_1,V_1$) and ($v_3,V_3$) differ as explained in the
Footnote~\ref{note1}. The symmetry can be restored if $b_{13}$ in
\eqref{vV1Cl30}-\eqref{sysACl30} is replaced by $-b_{31}$.},
\begin{align}
v_1&=\frac{b_1s+b_{23}S}{2(s^2+S^2)},
 &v_2&=\frac{b_{2}s-b_{13}S}{2(s^2+S^2)},
  &v_3&=\frac{b_3s+b_{12}S}{2(s^2+S^2)}, \label{vV1Cl30} \\
V_1&=\frac{b_{23}s-b_{1}S}{2(s^2+S^2)},
&V_2&=-\frac{b_{13}s+b_{2}S}{2(s^2+S^2)},
&V_3&=\frac{b_{12}s-b_{3}S}{2(s^2+S^2)}.\label{vV2Cl30}
\end{align}
The Eqs. \eqref{vV1Cl30}-\eqref{vV2Cl30} express the vectors $\bv$
and $\bV$ in terms of scalars $s$ and $S$, which are to be
determined from a pair of equations~\eqref{MVM1Cl30}. The
solution is valid when $s^2+S^2\neq 0$, i.e., when either $s\neq 0$
or $S\neq 0$, or both $s$ and $S$ are nonzero scalars. If these
conditions are not satisfied we have the case 3).  After
substitution of \eqref{vV1Cl30}--\eqref{vV2Cl30}, i.e.,
$(v_1,v_2,v_3)$ and $(V_1,V_2,V_3)$, into~\eqref{MVM1Cl30} we get
a system of two coupled algebraic equations for two unknowns $s$
and $S$,
\begin{equation}\begin{split}\label{sysACl30}
4(b_0-s^2+S^2)(s^2+S^2)^2=&+(b_1s+b_{23}S)^2+(b_2s-b_{13}S)^2+(b_3s+b_{12}S)^2\\
  &-(b_{23}s-b_1S)^2-(b_{13}s+b_2S)^2-(b_{12}s-b_3S)^2,\\
  2(b_{123}-2sS)(s^2+S^2)^2=&+(b_1s+b_{23}S)(b_{23}s-b_{1}S)-(b_{2}s-b_{13}S)\\
&\times(b_{13}s+b_{2}S)+(b_3s+b_{12}S)(b_{12}s-b_3S).
 \end{split}\end{equation}
If $s^2=S^2$ the solution of \eqref{sysACl30} requires some
compatibility condition to be satisfied since in this case we have
two equations and only one unknown. This situation will be
considered in the case 2) (see subsection~\ref{subsec:Cl30spec}). If $s$
and $S$ are independent, the case 1), it appears that
the system~\eqref{sysACl30} has exactly four solutions that can be
expressed in radicals. To this end we introduce new variables $t$
and $T$ and make the substitution
\begin{equation}\label{sysASCl30}
  s\,S = t,\qquad  \tfrac{1}{2}(-s^2 + S^2) = T,
\end{equation}
which reduces the system~\eqref{sysACl30} to
\begin{equation}\label{sysARCl30}
\begin{split}
&(b_{0}+4 T) (4 t-b_{123})-b_{I}=0,\\
&b_{S}-(b_{0}-b_{123}+4 T+4 t) (b_{0}+b_{123}+4 T-4 t)=0,
\end{split}\end{equation}
and where we have introduced a coordinate-free abbreviations
$b_{S}$ and $b_{I}$,
\begin{equation}\label{bSbICl30}
\begin{split}
&b_{S}= \langle \B \cliffordconjugate{\B}\rangle_0 =
b_{0}^2-b_{1}^2-b_{2}^2-b_{3}^2+b_{12}^2+b_{13}^2+b_{23}^2-b_{123}^2,\\
&b_{I}= \langle \B \cliffordconjugate{\B} I \rangle_0 = 2 b_{3}
b_{12}-2 b_{2} b_{13}+2 b_{1} b_{23}-2 b_{0} b_{123}.
\end{split}
\end{equation}
In \eqref{bSbICl30} the MV $\cliffordconjugate{\B}$ denotes the
Clifford conjugate of $\B$ and $\langle..\rangle_0$ is the scalar
filter. Note, that for remaining  algebras \cl{0}{3},
\cl{2}{1}, and \cl{2}{1} the signs of individual terms inside
$b_S$ and $b_I $ all are different. As we shall see below the
square  roots, in fact, are predetermined by four real
coefficients only, namely, the  scalar $b_0$ and
pseudoscalar~$b_{123}$ of MV $\B$, and signs in combination
of remaining coefficients in $b_{S}$ and~$b_{I}$
in~\eqref{bSbICl30}.


The substitution~\eqref{sysASCl30} and the resulting system of
equations~\eqref{sysARCl30} are of degree $\le 4$. The initial
system \eqref{sysACl30}, thus, can be explicitly solved in
radicals for CAs of  dimensions $n\le 3$. In particular, two real
solutions of \eqref{sysASCl30} have the form
\begin{equation}\label{solACl30}
\begin{split}
  &s_{1,2} =\pm\sqrt{-T+\sqrt{T^2+ t^2}},\qquad S_{1,2}
=\pm\frac{t}{\sqrt{-T+\sqrt{T^2+ t^2}}}\,.
\end{split}\end{equation}
  In~\eqref{solACl30}, the signs in pairs $(s_i,S_i)$ must be the
same (plus or minus). The denominator in $S_{1,2}$ becomes zero if
  $s=S=0$, i.e., in the case 3). The remaining two solutions of
\eqref{sysASCl30}, which can be obtained from \eqref{solACl30} by
the substitution $\sqrt{T^2+ t^2}\to -\sqrt{T^2+ t^2}$, are
complex valued, due to the inequality $\sqrt{T^2+ t^2}\ge T$, and
  therefore were rejected.

The two real valued solutions of Eq.~\eqref{sysARCl30} are
\begin{equation}\label{solBCl30}
\begin{split}
  &t_{1,2}=\frac{1}{4} \Bigl(b_{123} \pm \frac{1}{\sqrt{2}}\sqrt{-b_{S}+\sqrt{D}}\Bigr),\
  T_{1,2}=\frac{1}{4}\Bigl(\frac{\pm b_{I}}{\sqrt{2}\sqrt{-b_{S}+\sqrt{D}}}-
  b_{0}\Bigr),
\end{split}
\end{equation}
since for \cl{3}{0} algebra the determinant norm
$D=b_{S}^2+b_{I}^2$ of $\B$ is always positively definite and
$D\ge b_{S}$~\cite{Acus2018}. Again, we should take the same sign
for both $t_i$ and $T_i$. The two complex valued solutions of
\eqref{sysARCl30}, which can be obtained from \eqref{solBCl30} by
substitution $\sqrt{D}\to -\sqrt{D}$, must be rejected. The
  denominator of $T_{1,2}$ in \eqref{solBCl30} turns into zero when
$b_{S}=\sqrt{D}$, i.e., when $b_{I}=0$. In this case the solutions
of  \eqref{sysARCl30} are
\begin{equation}\label{solBd01Cl30}
\begin{split}
 t_{1,2}&=\begin{cases}
   \frac{1}{4} \bigl(\pm\sqrt{-b_{S}}+b_{123}\bigr),& b_{S}<0,\\
  b_{123}/{4} , &b_{S}>0.
 \end{cases}\quad
  T_{1,2}= \begin{cases}  -b_{0}/{4}, &b_{S}<0,\\
    \frac{1}{4} \bigl(\pm\sqrt{b_{S}}-b_{0}\bigr),& b_{S}>0.
  \end{cases}
\end{split}
\end{equation}
To summarize, beginning from \eqref{solBCl30} or
\eqref{solBd01Cl30} and then going to \eqref{solACl30}, and
finally to formulas \eqref{vV1Cl30}-\eqref{vV2Cl30}, one obtains
four explicit real solutions which completely determine the square
root $\sqrt{\B}=\A=s+\bv+(S+\bV)I$ of the generic MV $\B$ of real
$\cl{3}{0}$ algebra in terms of radicals.

\subsection{The special case $s^2 = S^2$\label{subsec:Cl30spec}}
The special case includes the subcases 1) $s = S\neq 0$, 2) $s =
-S\neq 0$ and 3)~$s=S=0$.
\subsubsection{The subcase $s=S\neq 0$}
  In this case the Eq.~\eqref{sysACl30} reduces to system
\begin{equation}\begin{split}\label{sysAspec1Cl30}
  (4 s^2 b_0-b_3 b_{12}+b_2 b_{13}-b_1 b_{23})/s=0,\\
  (16 s^4-b_1^2-b_2^2-b_3^2+b_{12}^2+b_{13}^2+b_{23}^2-8 s^2 b_{123})/s=0,
 \end{split}\end{equation}
where division by $s$ implies that the case $s=0$ should be
investigated separately. If the shortcuts $b_I$ and $b_S$ are
introduced and the requirement $s\neq 0$  is taken into account
the Eqs.~\eqref{sysAspec1Cl30} become
\begin{equation}\label{sysAspec30}
  b_I=-2b_0 (-4 s^2+ b_{123}),\qquad
  b_S=b_0^2-(-4 s^2+b_{123})^2.
\end{equation}
These equations are mutually compatible when the coefficients of
MV~$\B$ satisfy the condition $b_{I}^2+4 b_0^2(b_S- b_0^2)=0$. In
this case the solution is
\begin{equation}\begin{split}\label{solSspec1Cl30}
s_{1,2} =&\begin{cases}
\pm\tfrac12\sqrt{b_{123}+\dfrac{b_I}{2 b_0}} & \textrm{if}\quad b_0\neq  0,\\
\pm\tfrac12\sqrt{b_{123}\pm\sqrt{-b_S}} & \textrm{if}\quad b_0 =  0,
  \end{cases}
\end{split}\end{equation}
where all expressions inside square roots are assumed to be
positive. The expressions  \eqref{solSspec1Cl30} and
\eqref{vV1Cl30}-\eqref{vV2Cl30} yield the final answer for this
subcase.

\subsubsection{The subcase $s = -S\neq 0$}
When $s=-S$ we proceed in the same way.  Now we have to solve
the equations
\begin{equation}\label{sysAspec}
 b_I=-2b_0 (4 s^2+ b_{123}),\qquad
 b_S=b_0^2-(4 s^2+b_{123})^2.
\end{equation}
The equations have a solution if coefficients of $\B$ satisfy the
same condition $b_{I}^2+4 b_0^2(b_S- b_0^2)=0$. The solution then
is
\begin{equation}\begin{split}\label{solSspec2}
s_{1,2} =&\begin{cases}
\pm\tfrac12\sqrt{-b_{123}-\dfrac{b_I}{2 b_0}} & \textrm{if}\quad b_0\neq  0,\\
\pm\tfrac12\sqrt{-b_{123}\pm\sqrt{-b_S}} & \textrm{if}\quad b_0 =  0,
  \end{cases}
\end{split}\end{equation}
where again all expressions inside square roots must  be positive.
The Eqs.~\eqref{solSspec2} and \eqref{vV1Cl30}-\eqref{vV2Cl30}
yield the square root of $\B$ in this special case. The
expressions \eqref{solSspec2} and \eqref{solSspec1Cl30} are
complementary, i.e., for a given $\B$ they cannot
 occur simultaneously.

\subsubsection{The subcase $s = S = 0$}
This special case is interesting because the condition $s = S = 0$
implies that the number of square roots $\sqrt{\B}$ may be
infinite. Indeed, when $s = S = 0$ the Eqs.
\eqref{MVM2Cl30}-\eqref{MVM4Cl30} are compatible only if the
vector coefficients $b_1, b_2, b_3$ and bivector coefficients
$b_{12}, b_{13}, b_{23}$ all are equal to zero. Then, the
Eq.~\eqref{MVM1Cl30} reduces to
\begin{equation}\label{solSspec3Cl30}
b_0=\bv^2-\bV^2,\quad b_{123}=2(\bv\cdot\bV).
\end{equation}
Since, in general, we have $3+3=6$ unknowns which have to satisfy
only two equations we are left with four real parameters which can
be chosen at will (see Example~1). The solution set, therefore,
makes a four dimensional (or smaller) manifold as long as manifold
can ensure real-valued MV coefficients. Here we do not investigate
the range of values of the free parameters. It should be noticed
that  $\B = b_0+b_{123}\e{123}$ belongs to a center of \cl{3}{0}
algebra. Two Eqs.~\eqref{solSspec3Cl30} have a clear geometric
interpretation. Indeed, imagine two concentric spheres determined
by end points of vectors $\bv$ and $\bV$.  The coefficient $b_0$
then determines difference of lengths of $|\bv|^2$ and $|\bV|^2$,
while the coefficient at pseudoscalar $b_{123}$ controls an angle
between the vectors $\bv$ and~$\bV$. Before considering
non-Euclidean algebras we provide some examples.
%
\subsection{Examples for $\cl{3}{0}$}
\subsubsection*{Example~1. The case $s\ne S$.}\label{exmp1}
The square root of $\B=\e{1}-2 \e{23}$. The coefficients in this
case are $b_1=1$ and $b_{23}=-2$, and all remaining are equal to
zero. Then from \eqref {bSbICl30} we have that  $b_I=-4$ and
$b_S=3$. The expression \eqref{solBCl30} gives
$t_{1,2}=(\tfrac{1}{4},-\tfrac{1}{4})$ and
$T_{1,2}=(-\tfrac{1}{2},\tfrac{1}{2})$.  Finally using
\eqref{solACl30} we find the real solutions for $s$ and $S$,
\begin{equation}\begin{split}\label{ex1}
  &\bigl(s_{1,2}=\mp\tfrac{1}{2}c_1,\ S_{1,2}=\pm\tfrac{1}{2}c_2\bigr)\quad \text{and}\quad
  \Bigl(s_{3,4}=\pm\tfrac{1}{2}c_2,\ S_{3,4}=\pm\tfrac{1}{2}c_1\Bigr),
\end{split}\end{equation}
where $c_1=\sqrt{-2+\sqrt{5}}$ and  $c_2=\sqrt{2+\sqrt{5}}$. The
MV, therefore, is regular and we can omit the test of special
cases $s^2=S^2$. Using \eqref{vV1Cl30}-\eqref{vV2Cl30} then we
have the following four sets of non-zero coefficients:
\begin{equation}\label{ex1aCl30}\begin{array}{llll}
\bigl(s_1=-\tfrac{1}{2}c_1,& S_1=\tfrac{1}{2}c_2,& v_1=-\tfrac{1}{2}c_2,& V_1=-\tfrac{1}{2}c_1\bigr),\\
\bigl(s_2=\tfrac{1}{2}c_1,& S_2=-\tfrac{1}{2}c_2,& v_1=\tfrac{1}{2}c_2,& V_1=\tfrac{1}{2}c_1\bigr),\\
\bigl(s_3=\tfrac{1}{2}c_2,& S_3=\frac{1}{2}c_1,& v_1=\tfrac{1}{2}c_2,& V_1=-\tfrac{1}{2}c_1\bigr),\\
\bigl(s_4=-\tfrac{1}{2}c_2,& S_4=-\tfrac{1}{2}c_1,&
v_1=-\tfrac{1}{2}c_1,& V_1=\tfrac{1}{2}c_2\bigr).
\end{array}\end{equation}
The remaining coefficients are equal to zero, $v_2=v_3=V_2=V_3=0$.
Finally, inserting the coefficients~\eqref{ex1aCl30} into
\eqref{mvA30} we find all four roots,
\begin{equation}\begin{split}\label{ex1fa}
  &\sqrt{\B_{1,2}}=\A_{1,2} = \mp\frac12 c_2 \bigl(-2+\sqrt{5}+\e{1}+(-2+\sqrt{5}) \e{2 3}-\e{1 2 3}\bigr),\\
  &\sqrt{\B_{3,4}}=\A_{3,4}= \pm \frac12 c_1 \bigl(2+\sqrt{5}+\e{1}-(2+\sqrt{5}) \e{2 3}+\e{1 2
  3}\bigr),
\end{split}\end{equation}
the squares of which  give the initial MV $\B=\e{1}-2 \e{23}$.

\subsubsection*{Example~2. The case $s=S$.\label{exmp2}}
The square root of $\B=-1+\e{3}-\e{12}+\frac12\e{123}$. We have
$b_0=-1$, $b_{123}=\tfrac12$, $b_I=-1$, $b_S=\tfrac34$. From
\eqref{solACl30} and \eqref{solBCl30} then we find
\begin{equation}\label{ex2}
\bigl(s_{1,2}=\pm\tfrac{1}{2},\quad S_{1,2}=\pm\tfrac{1}{2}\bigr).
\end{equation}
Since $s=S$ and the condition $b_{I}^2+4 b_0^2(b_S- b_0^2)=0$ is
satisfied, from~\eqref{solSspec1Cl30} we recover once again
$s_{1,2}=\pm\frac12$, whereas the Eq.~\eqref{solSspec2} has no
real solution. Then equations \eqref{vV1Cl30}-\eqref{vV2Cl30}
yield
\begin{equation}\begin{array}{llll}\label{ex2a}
\bigl(s_1=-\tfrac{1}{2},&\ v_1=v_2=v_3=0,&\ V_1=V_2=0,&\ V_3=1\bigr),\\
\bigl(s_2=\tfrac{1}{2},&\ v_1=v_2=v_3=0,&\ V_1=V_2=0,&\ V_3=-1\bigr).\\
\end{array}\end{equation}
The answer, therefore, is
\begin{equation}\begin{split}\label{ex2fa}
  &\sqrt{\B_{1,2}} =\A_{1,2}= \pm \bigl( \tfrac12 \e{3} + \tfrac12 \e{12}-\e{1 2 3}\bigr),\\
  &\sqrt{\B_{3,4}}= \A_{3,4}= \pm \bigl(-\tfrac12 +\e{12}-\tfrac12 \e{1 2 3}\bigr).\\
\end{split}\end{equation}


\subsubsection*{Example~3. The case $s=S=0$ and $s\ne S\ne 0$.}\label{ex3Cl30}
The square root of MV $\B=-1+\e{123}$ that represents the
algebra center.
 We have $b_0=-1$, $b_{123}=1$ and $b_I=2$, $b_S=0$. Then
from expressions  \eqref{solBCl30} and \eqref{solACl30} we obtain
\begin{equation}\label{ex3}\begin{split}
&\bigl(s_{1,2}=\pm c_1,\ S_{1,2}=\pm c_2\bigr)\quad
\text{and}\quad  \bigl(s_{3}=0,\ S_{3}=0\bigr),
\end{split}\end{equation}
where  $c_1=\sqrt{-1/2+1/\sqrt{2}}$ and
$c_2=\sqrt{1/2+1/\sqrt{2}}$. The cases $s_{1,2}\neq S_{1,2}$ in
\eqref{ex3} can be computed as in the Example~1 above.
The square roots are $\sqrt{\B_{1,2}} = \pm (c_1+c_2 \e{123})$.

The case $(s_{3}=0,\ S_{3}=0)$ in \eqref{ex3} is special. Since
all coefficients ($b_1, b_2, b_3$) and ($b_{12}, b_{13}, b_{23}$)
are zeroes the solution set is not empty. Indeed, as seen  from
\eqref{solSspec3Cl30} the system can  be solved for an arbitrary
pair of coefficients ($v_1,v_2,v_3,V_1,V_2,V_3$), for example, for
($v_1,V_1$). After substitution of ($v_1,V_1$) into~\eqref{mvA30}
we obtain MV with four free parameters
$\sqrt{\B_{3}}=f_1(v_2,v_3,V_2,V_3)\e{1}+v_2\e{2}+v_3\e{3}+f_2(v_2,v_3,V_2,V_3)\e{23}-V_2\e{13}+V_3\e{12}$,
where $v_1=f_1(v_2,v_3,V_2,V_3)$ and $V_1=f_2(v_2,v_3,V_2,V_3)$
denote explicit solutions of system~\eqref{solSspec3Cl30},
\begin{equation}\label{exampleCl30free}
\begin{split}
&v_{1}=\mp \frac{c_1}{\sqrt{2}},\qquad
  V_{1}=\pm \frac{1}{c_1} \frac{-b_{123}+2 (v_{2} V_{2}+v_{3} V_{3})}{\sqrt{2}},\qquad \textrm{with}\quad \\[3pt]
&c_1=\Bigl(\pm\sqrt{\left(b_{0}-v_{2}^2-v_{3}^2+V_{2}^2+V_{3}^2\right)^2+
   (b_{123}-2 (v_{2} V_{2}+v_{3} V_{3}))^2}\\
&\phantom{\frac{1}{\sqrt{2}}b_{0}-v_{2}^2-v_{3}^2  (b_{123}-2
(v_{2} V_{2}+v_{3} V_{3}))x}
 +b_{0}-v_{2}^2-v_{3}^2+V_{2}^2+V_{3}^2\Bigr)^{\frac{1}{2}}
.
\end{split}
\end{equation}
If instead of $\B=-1+\e{123}$ we would attempt to find the square
root of MV that does not belong to the center, for example
$\B=\e{1}+\e{12}$ (which represents physically important case of
linearly polarized electromagnetic wave in $\cl{3}{0}$), we would
end with an empty solution set. Indeed, in this case we have
$\bigl(s_1=0,\ S_1=0\bigr)$ and $b_0=b_{123}=b_I=b_S=0$,
therefore, the compatibility condition $b_{I}^2+4 b_0^2(b_S-
b_0^2)=0$ is satisfied. The Eqs.~\eqref{vV1Cl30}-\eqref{vV2Cl30},
however, yields contradiction $1=0$. This is not surprising, since
for this particular MV we have $(\e{1}+\e{12})^2=0$.


\subsubsection*{Example~4. The case of quaternion.}
Since quaternions are isomorphic to even subalgebra
$\cl{3}{0}^{+}$ with  elements $\{1,\e{12},\e{23}, \e{13}\}$, the
provided formulas allow to find the square root of a quaternion as
well. For example, taking into account that $\mathbf{i} =\e{12}$,
$\mathbf{j}=-\e{13}$, $\mathbf{k}=\e{23}$ let us compute the
square root of
$\B=1+\e{12}-\e{13}+\e{23}=1+\mathbf{i}+\mathbf{j}+\mathbf{k}$. In
this case we have  $b_0=1$, $b_{123}=0$ and $b_I=0$, $b_S=4$.
Beginning  from \eqref{solBd01Cl30} and using Eq.~\eqref{solACl30}
we find that this is a regular case with four different
coefficients
\begin{equation}\begin{split}\label{ex4}
  &\bigl(s_{1,2}=0,\quad S_{1,2}=\pm 1/\sqrt{2}\bigr)\quad \text{and}\quad
  \Bigl(s_{3,4}=\pm \sqrt{3/2},\quad S_{3,4}=0\Bigr).
\end{split}\end{equation}
Using \eqref{vV1Cl30}-\eqref{vV2Cl30} and \eqref{mvA30} we can write the answer:
\begin{equation}\begin{split}\label{ex4fa}
  &\sqrt{\B_{1,2}} =\A_{1,2}= \pm (\e{1} +\e{2}+\e{3}+\e{1 2 3})/\sqrt{2},\\
  &\sqrt{\B_{3,4}}= \A_{3,4}= \pm(3 +\e{12}-\e{13}+\e{23})/\sqrt{6}=\pm(3 +\mathbf{i}+\mathbf{j}+\mathbf{k})/\sqrt{6}.
\end{split}\end{equation}
The squares  of all four roots gives the initial MV. It should be
noticed that the roots $\sqrt{\B_{3,4}}$ have remained in the even
algebra, whereas  the roots $\sqrt{\B_{1,2}}$ have not. The source
of this difference is related with the initial assumption that the
roots belong to  \cl{3}{0} algebra rather than to quaternionic or
isomorphic \cl{0}{2} algebra.


\section{\cl{0}{3} algebra}\label{sec:Cl03}
A similar calculation route allows to write down explicit square
root formulas for \cl{0}{3} algebra as well. Using the same
notation~\eqref{mvA30} for $\A$ and $\B$ and equating coefficients
at the same basis elements in $\A^2=\B$ we obtain the
equations,\footnote{The reader has noticed that formulas
 have similar structure and only differ in signs of
some constituent terms. For easier reading and application, all
formulas below -- including that for mixed algebras -- are written
explicitly without introducing a large number of sign epsilons
such as $\varepsilon_{\pm}=\pm 1$. The appearance of different
signs in structurally similar expressions, in fact, brings in
different conditions for existence of real roots in distinct
algebras.}
\begin{align}
b_0&=s^2+S^2+\bv^2+\bV^2, &b_{123}&=2(sS+\bv\cdot\bV)\label{MVM10Cl03}, \\
b_1&=2(s v_1+S V_1), &b_{23}&=-2(s V_1+S v_1)\label{MVM20Cl03}, \\
b_2&=2(s v_2+S V_2), &b_{13}&=2(s V_2+S v_2)\label{MVM30Cl03}, \\
b_3&=2(s v_3+S V_3), &b_{12}&=-2(s V_3+S v_3)\label{MVM40Cl03},
\end{align}
where now $\bv^2=-v_1^2-v_2^2-v_3^2$ and $\bv\cdot\bV=-v_1 V_1- v_2 V_2- v_3 V_3$.


\subsection{The generic case, $s^2\neq S^2$}
The solution of  Eqs.~\eqref{MVM20Cl03}-\eqref{MVM40Cl03} is
\begin{align}
v_1&=\frac{b_1s+b_{23}S}{2(s^2-S^2)}, & v_2&=-\frac{b_{2}s- b_{13}S}{2(s^2-S^2)},&  v_3&=\frac{b_3s+b_{12}S}{2(s^2-S^2)},\label{v123Cl03}\\
V_1&=-\frac{b_{23}s+b_{1}S}{2(s^2-S^2)},
&V_2&=\frac{b_{13}s-b_{2}S}{2(s^2-S^2)},
&V_3&=-\frac{b_{12}s+b_{3}S}{2(s^2-S^2)}.\label{V123Cl03}
\end{align}
The solution is valid when $s^2-S^2\neq 0$ and corresponds to
generic case considered in this subsection. We substitute
\eqref{v123Cl03} and \eqref{V123Cl03} into \eqref{MVM10Cl03} and
obtain two coupled algebraic equations for two unknowns $s$ and
$S$,
\begin{equation}\label{sysA0Cl03}
\begin{split}
& b_{S}+4 s^2 (-6 S^2+b_0)+8 s S b_{123}=4 s^4+(-2 S^2+b_0)^2+b_{123}^2,\\
  & b_{I}=2 (2 (s^2+S^2)-b_0) (4 s S-b_{123}),
\end{split}\end{equation}
where we again have introduced the coordinate-free notation for
\begin{equation}\label{bSbICl03}
\begin{split}
&b_{S}= \langle \B \cliffordconjugate{\B}\rangle_0 =
b_{0}^2+b_{1}^2+b_{2}^2+b_{3}^2+ b_{12}^2+b_{13}^2+b_{23}^2+b_{123}^2,\\
&b_{I}= \langle \B
\cliffordconjugate{\B} I \rangle_0 = -2 b_{3} b_{12}+2 b_{2} b_{13}-2 b_{1} b_{23}+2 b_{0} b_{123}.
\end{split}
\end{equation}
(Note change of signs compared to \cl{3}{0} case). The determinant norm in \cl{0}{3} is expressed as a difference,
$D=b_S^2-b_I^2$, and is always positive $D>0$.

To reduce the degree of the above equations we make the
substitution
\begin{equation}\label{sysAS0Cl03}
  s\,S = t;\qquad  \tfrac{1}{2}(s^2 + S^2) = T,
\end{equation}
which transforms\footnote{To eliminate $s$ and $S$ {\it
Mathematica} commands \textbf{Eliminate[~]} and
  \textbf{GroebnerBasis[~]} can be used. These commands allow to rewrite the
initial system of equations in a number of equivalent forms.} the
  system~\eqref{sysA0Cl03} to simpler one,
\begin{equation}\label{sysAR0Cl03}
 b_{S}=(4 T-b_0)^2+(4 t-b_{123})^2,\qquad
  b_{I}=2 (4 T-b_0) (4 t-b_{123}).
\end{equation}
The solution of~\eqref{sysAR0Cl03} when $(b_{S}\pm\sqrt{D})
> 0$ is
\begin{equation}\label{solB0Cl03}
\begin{split}
  &t_{1,2}=\frac{1}{4} \Bigl(b_{123} \pm \frac{1}{\sqrt{2}}\sqrt{b_{S}-\sqrt{D}}\Bigr),\quad
  T_{1,2}=\frac{1}{4}\Bigl(b_{0}\pm\frac{ b_{I}}{\sqrt{2}\sqrt{b_{S}-\sqrt{D}}}\Bigr).
\end{split}
\end{equation}
We have to take either plus or minus signs in $t_i$ and
$T_i$ (two possibilities). The remaining two solutions of
\eqref{sysAR0Cl03}, which were obtained
from \eqref{solB0Cl03} by the replacement $-\sqrt{D}\to +\sqrt{D}$
yielded a complex valued expression for $T\pm\sqrt{T^2-
t^2}$ (see Eq.~\eqref{solA0Cl03} below) and were dismissed
in advance.

When  $b_{I}=0$ or, equivalently, $b_{S}\pm\sqrt{D} = 0$ the
solution of \eqref{sysAR0Cl03} is
\begin{equation}\label{solB0d01aCl03}
  (t_{1,2},\ T_{1,2})= \bigl(\frac{1}{4}b_{123},\ \frac{1}{4} (\pm\sqrt{b_{S}}+b_{0})\bigr),
\end{equation}
which, alternatively, can be also written in a form
\begin{equation}\label{solB0d01bCl03}
  (t_{1,2},\ T_{1,2})= \bigl(\frac{1}{4}b_0,\ \frac{1}{4} (\pm\sqrt{b_{S}}+b_{123})\bigr).
\end{equation}
Once equations in~\eqref{solB0Cl03} are computed they can be
substituted back into solutions of~\eqref{sysAS0Cl03}
\begin{equation}\label{solA0Cl03}
s_{1,2,3,4} =\pm\sqrt{T\pm\sqrt{T^2- t^2}},\;\, S_{1,2,3,4}
  =\pm\frac{t}{\sqrt{T\pm\sqrt{T^2- t^2}}}, \textrm{ for }T\ge 0.
\end{equation}
We have to choose the same signs in the same positions for $s_{1,2,3,4}$
and $S_{1,2,4,5}$ (four possibilities). In the case when
$T\pm\sqrt{T^2- t^2}=0$ (i.e., for $t=0$) the solutions of
equations in~\eqref{sysAS0Cl03} are the sets
\begin{equation}\label{solA0d01Cl03}
\begin{split}
& (s_{1,2}= S_{3,4} = \pm\sqrt{2T}, \qquad S_{1,2}= s_{3,4}  = 0).
\end{split}\end{equation}
Thus, starting  from pairs $(t_1,T_1)$ and $(t_2,T_2)$ in
Eq.~\eqref{solB0Cl03} (or \eqref{solB0d01aCl03}, equivalently
\eqref{solB0d01bCl03}) and then passing to \eqref{solA0Cl03} (or
\eqref{solA0d01Cl03}), and finally to formulas \eqref{v123Cl03},
\eqref{V123Cl03} and \eqref{mvA30} one obtains explicit real
solutions which completely determine the square root
$\sqrt{\B}=\A=s+\bv+(S+\bV)I$ of the generic MV $\B$ of real
$\cl{0}{3}$ algebra in terms of radicals. It appears that only
four\footnote{The condition $T\ge 0$ in \eqref{solA0Cl03} selects
only one sign from \eqref{solB0Cl03}.} real solutions at most are
possible in this algebra too, since other choices of signs in
\eqref{solB0Cl03} and \eqref{solA0Cl03} yield negative expressions
inside square roots. Now we consider special case.

\subsection{The special case $s^2 = S^2$}

We have the subcases  1) $s = S\neq 0$,  2)  $s = -S\neq 0$ and 3) $s=S=0$.
\subsubsection{The subcase $s=S\neq 0$}
The system of  Eqs.~\eqref{MVM20Cl03}--\eqref{MVM40Cl03} has a
special solution,
\begin{equation}\begin{split}\label{vVSpecialCl03}
v_{1}= \frac{b_{1}}{2 s}-V_{1},\quad v_{2}= \frac{b_{2}}{2 s}-V_{2},\quad v_{3}= \frac{b_{3}}{2 s}-V_{3},
\end{split}\end{equation}
iff the MV $\B$ coefficients satisfy $b_1 =-b_{23}$, $b_2 =
b_{13}$,  $b_3= - b_{12}$. In \eqref{vVSpecialCl03} we have
expressed $v_i$ in terms of $V_i$. Appearance of $s$ in the
dominators of \eqref{vVSpecialCl03} implies that the case $s=S=0$
should again be investigated separately. After substituting the
solution \eqref{vVSpecialCl03} into \eqref{MVM10Cl03} and taking
into account the mentioned conditions $b_1 =- b_{23}, b_2 =
b_{13}, b_3= - b_{12}$ we get two equations,
\begin{equation}\begin{split}\label{sysAspecCl03}
  &-\frac{b_{1}^2+b_{2}^2+b_{3}^2}{4 s^2}+\frac{b_{1} V_{1}+b_{2}
  V_{2}+b_{3} V_{3}}{s}-b_{0}+2 s^2+2 (\bV\cdot \bV)=0,\\
&-\frac{b_{1} V_{1}+b_{2}  V_{2}+b_{3} V_{3}}{s}-b_{123}+2 s^2-2 (\bV\cdot \bV)=0,
\end{split}\end{equation}
which should be  kept mutually compatible. To this end we subtract
and add the above equations,
\begin{equation}\begin{split}\label{sysAspecSubsAddCl03}
  &-\frac{4 s^2 \left(b_{0}-b_{123}-4 (\bV\cdot \bV)\right)-8 s (b_{1} V_{1}+b_{2} V_{2}+b_{3} V_{3})+b_{1}^2+b_{2}^2+b_{3}^2}{4 s^2}=0,\\
  &  -\frac{4 s^2 \left(b_{0}+b_{123}-4 s^2\right)+b_{1}^2+b_{2}^2+b_{3}^2}{4 s^2}=0.
\end{split}\end{equation}
Then using the explicit form of the relation $b_0^2 + 2 (b_1^2 + b_2^2
+ b_3^2) + b_{123}^2 = b_S=\langle \B
  \cliffordconjugate{\B}\rangle_0$, where we have taken into account the
conditions $b_1 =- b_{23},  b_2 = b_{13},  b_3= - b_{12}$, we
express $b_{1}^2+b_{2}^2+b_{3}^2$ from the second equation of
\eqref{sysAspecSubsAddCl03}, $b_{1}^2+b_{2}^2+b_{3}^2 =\frac12
(b_S-b_{0}^2-b_{123}^2)$, and substitute it into the first of
equations in~\eqref{sysAspecSubsAddCl03}. The result is a
quadratic equation for $V_i$. After solving it, for example, with
respect to $V_1$ we express it in terms of now arbitrary free
parameters $V_2$ and $V_3$,
\begin{equation}\begin{split}\label{AspecCl03}
  V_{1}=& \frac{\sqrt{2}}{8 s}\Bigl(\sqrt{2} b_{1}\pm \bigl(-8 s^2 \left(b_{0}-b_{123}+4 \left(V_{2}^2+V_{3}^2\right)\right)+\\
  &\phantom{\bigl(-8 s^2 }
  16 s (b_{2} V_{2}+b_{3} V_{3})+b_{0}^2+2 b_{1}^2+b_{123}^2-b_{S}\bigr)^{1/2}\Bigr),
\end{split}\end{equation}
which ensures the compatibility of the system
\eqref{sysAspecCl03}. Therefore we can restrict further analysis
to a single simplest equation, which we choose to be the second
equation of \eqref{sysAspecSubsAddCl03}. Using the shortcut $b_S$
it can be cast in the form,
  \begin{equation}\begin{split}\label{sysAspec2Cl03}
b_{S}=-8 b_{0} s^2-8 b_{123} s^2+b_{0}^2+b_{123}^2+32 s^4 .
 \end{split}\end{equation}
After solving it with respect to $s$ we have
\begin{equation}\begin{split}\label{solSspec1Cl03}
s_{1,2} =& \pm\tfrac{1}{2 \sqrt{2}}\sqrt{\sqrt{2
b_{S}-(b_{0}-b_{123})^2}+b_{0}+b_{123}}\,,
\end{split}\end{equation}
where all expressions inside square roots are assumed to be
positive (we again have removed complex valued solution). The
expressions \eqref{solSspec1Cl03}, \eqref{AspecCl03} and
\eqref{vVSpecialCl03} substituted into \eqref{mvA30} yield the
final answer for this special case provided that the conditions
$b_1 =- b_{23},  b_2 = b_{13},  b_3= - b_{12}$, which can be
shortly written also as $b_S-(b_0-b_{123})^2=b_I$, are satisfied. The
solution set contains two free parameters (here $V_2$ and $V_3$)
and, therefore, makes two (or smaller) dimensional manifold in
the parameter space.

\subsubsection{The subcase $s=-S\neq 0$}
Performing exactly the same analysis we obtain that solution in
this case exists when the conditions $b_1 = b_{23}$, $b_2 =
-b_{13}$ and $b_3= b_{12}$, which can be shortly written as
$-b_S+(b_0+b_{123})^2=b_I$, are satisfied. Then we can express
$v_i$ in terms of $V_i$
\begin{equation}\begin{split}\label{vVSpecial2Cl03}
v_{1}= \frac{b_{1}}{2 s}+V_{1},\quad v_{2}= \frac{b_{2}}{2 s}+V_{2},\quad v_{3}= \frac{b_{3}}{2 s}+V_{3}.
\end{split}\end{equation}
In a similar way $V_1$ can  be expressed in terms of $V_2$ and
$V_3$
\begin{equation}\begin{split}\label{Aspec2Cl03}
  V_{1}=& \frac{\sqrt{2}}{8 s}\Bigl(-\sqrt{2} b_{1}\pm \big(-8 s^2 \big(b_{0}+b_{123}+4(V_{2}^2+V_{3}^2)\big)\\
  &\phantom{\big(-8 s^2 }
  -16 s (b_{2} V_{2}+b_{3} V_{3})+b_{0}^2+2 b_{1}^2+b_{123}^2-b_{S}\big)^{1/2}\Bigr),
\end{split}\end{equation}
whereas for $s$ we have two real solutions
\begin{equation}\begin{split}\label{solSspec1NegCl03}
s_{1,2} =& \pm\tfrac{1}{2 \sqrt{2}}\sqrt{\sqrt{2
b_{S}-(b_{0}+b_{123})^2}+b_{0}-b_{123}}\,,
\end{split}\end{equation}
The above expressions substituted into \eqref{mvA30} again yield
the final MV provided the conditions $b_1 = b_{23},  b_2 =
-b_{13}, b_3= b_{12}$ are satisfied.

\subsubsection{The subcase $s = S = 0$}
The analysis of this special subcase is very similar to that in
$\cl{3}{0}$. The Eqs.~\eqref{MVM20Cl03}-\eqref{MVM40Cl03} are
compatible if  the  coefficients at vector  $b_1, b_2, b_3$ and
bivector $b_{12}, b_{13}, b_{23}$  are equated to zero. Then we
are left with Eqs.~\eqref{MVM10Cl03} which assume the following
form:
\begin{equation}\label{solSspec3Cl03}
b_0=\bv^2+\bV^2,\quad b_{123}=2(\bv\cdot\bV).
\end{equation}
From these equations follows that  four parameters are left
unspecified. For example, we can again solve
Eq.~\eqref{solSspec3Cl03} with respect to pair $(v_1,V_1)$,
\begin{equation}\begin{split}
&v_{1}=\mp \frac{c_1}{\sqrt{2}},\qquad
  V_{1}=\pm \frac{1}{c_1} \frac{b_{123}+2 (v_{2} V_{2}+v_{3} V_{3})}{\sqrt{2}},\qquad \textrm{where}\quad \\[3pt]
&c_1=\Bigl(\pm\sqrt{\left(b_{0}+v_{2}^2+v_{3}^2+V_{2}^2+V_{3}^2\right)^2-
   (b_{123}+2 (v_{2} V_{2}+v_{3} V_{3}))^2}\\
&\phantom{\frac{1}{\sqrt{2}}b_{0}-v_{2}^2-v_{3}^2  (b_{123}-2
(v_{2} V_{2}+v_{3} V_{3}))x}
-b_{0}-v_{2}^2-v_{3}^2-V_{2}^2-V_{3}^2\Bigr)^{\frac{1}{2}}.
\end{split}
\end{equation}
The pairs  $(v_2,V_2)$ and $(v_3,V_3)$ can be interpreted as
free parameters that generate a continuum of roots in a four
parameter space.

\subsection{Examples for $\cl{0}{3}$}
\subsubsection*{Example~5. The case $s\ne S$.}
Square root of the same MV $\B=\e{1}-2 \e{23}$ as in Example~1.
The nonzero coefficients are $b_1=1$, $b_{12}=-2$ and the
shortcuts in~\eqref{bSbICl03} acquire values $b_I=4$ and $b_S=5$.
The expression \eqref{solB0Cl03} gives $(T_{1},t_{1})=
(\tfrac{1}{2},\tfrac{1}{4})$ and $(T_{2},t_{2})=
(-\tfrac{1}{2},-\tfrac{1}{4})$. Because $T_1>0$ we take
$(T_{1},t_{1})$. Then from \eqref{solA0Cl03} we find four values
of $(s_i,S_i)$ and then from \eqref{v123Cl03}, \eqref{V123Cl03}
the coefficients $v_i$ and $V_i$,
\begin{equation}\label{ex1aCl03}
\begin{array}{lll}
\big(s_1=V_1=\tfrac{1}{4}d_3,&S_1=v_1=-\tfrac{1}{2}d_2,&v_2=v_3=V_2=V_3=0\big),\\
\big(s_2=V_1=\tfrac{1}{2}d_1,&S_2= v_1=\tfrac{1}{2}d_2,&v_2=v_3=V_2=V_3=0\big),\\
\big(s_3=V_1=-\tfrac{1}{2}  d_2,&S_3=-\tfrac{1}{2 d_2},\ v_1=\tfrac{1}{4}d_3,&v_2=v_3=V_2=V_3=0\big),\\
\big(s_4= V_1=\tfrac{1}{2}d_2,&S_4=\tfrac{1}{2d_2},\ v_1=\tfrac{1}{2}d_1,& v_2=v_3=V_2=V_3=0\big).\\
\end{array}
\end{equation}
where $d_1=\sqrt{2-\sqrt{3}}$,  $d_2=\sqrt{2+\sqrt{3}}$ and
$d_3=\sqrt{2}-\sqrt{6}$. Finally, we find four square roots as in
Example~1.
\begin{equation}\begin{split}\label{ex1AnsCl03}
 &\A_{1}=\sqrt{\B_{1}}=\tfrac{1}{2}\big(d_1+d_2 \e{1}-d_1 \e{23}+ d_2\e{1 2 3}\big),\\
 &\A_{2}=\sqrt{\B_{2}}=\tfrac{1}{2}\big(-d_2+ \tfrac{1}{2}d_3\e{1}+d_2\e{2 3} -\frac{\e{1 2 3}}{d_2}\big),\\
 &\A_{3}=\sqrt{\B_{3}}=\tfrac{1}{2}\big(d_2+ d_1\e{1}-d_2\e{2 3}+\frac{\e{123}}{d_2}\big),\\
 &\A_{4}=\sqrt{\B_{4}}=\tfrac{1}{2}\big(\tfrac{1}{2}d_3-d_2\e{1}-\tfrac{1}{2}d_3 \e{23} -d_2\e{1 2 3}\big).\\
\end{split}\end{equation}
Though symbolical forms look different, in fact $\A_{1}=-\A_{4}$
and $\A_{2}=-\A_{3}$. By this reason there are four
different roots with either plus or minus signs rather then eight
roots with plus/minus signs.


\subsubsection*{Example~6. The case $s= S\neq 0$.}
The square root of  $\B=-\e{3}+\e{12}+4\e{123}$. The shortcuts
have the values $b_I=4$ and $b_S=2$. The expression
\eqref{solB0Cl03} gives $(T_{1},t_{1})=
(\tfrac{c_1}{4},\tfrac{c_1}{4})$, where $c_1=(2+\sqrt{5})$. The
negative $T$ solution was omitted. For $(s_i, S_i)$ then we have
$(s_1,S_1)=(-\tfrac{\sqrt{c_1}}{2},-\tfrac{\sqrt{c_1}}{2})$ and
$(s_2,S_2)=(\tfrac{\sqrt{c_1}}{2},\tfrac{\sqrt{c_1}}{2})$. The
coefficients satisfy the relation $b_1 =- b_{23}$,  $b_2 =
b_{13}$,  $b_3= -b_{12}$, therefore, a special solution exists.
Then for $s=-\tfrac{\sqrt{c_1}}{2}$ we find $v_1=-V_1=\frac{1}{2}
\sqrt{\frac{1}{2} c_2} d_1$, $v_2=-V_2$, $v_3=\sqrt{c_2}-V_3$,
with $c_2=(-2+\sqrt{5})$, where $d_1=\sqrt{-8 c_1
\left(V_2^2+V_3^2-1\right)+8 \sqrt{c_1} V_3-2}$. For
$s=\tfrac{\sqrt{c_1}}{2}$ we have $V_1=-v_1=\frac{1}{2}
\sqrt{\frac{1}{2} c_2} d_2$, where $d_2=\sqrt{-8 c_1
\left(V_2^2+V_3^2-1\right)-8 \sqrt{c_1} V_3-2}$ and $v_2=-V_2$,
$v_3=-V_3-\sqrt{c_2}$. The coefficients $V_2$ and $V_3$ represent
free parameters. Finally, we have the roots 
\begin{equation}\begin{split}\label{ex2AnsCl03}
\A_{1,2}=&\sqrt{\B_{1,2}}=\mp\frac{\sqrt{c_1}}{2}\pm\frac{d_{1,2}}{2\sqrt{2 c_1}}\e{1}-V_2\e{2} + \big(\pm\frac{1}{\sqrt{c_1}}-V_3\big)\e{3} -V_3 \e{12}\\
 &\phantom{-\frac{\sqrt{c_1}}{2}+\frac{d_{1,2}}{2\sqrt{2 c_1}}\e{1}-V_2\e{2}}
+V_2\e{13}\pm\frac{d_{1,2}}{2\sqrt{2
c_1}}\e{23}\mp\frac{\sqrt{c_1}}{2}\e{123}
\end{split}\end{equation}
Taking, for example, $V_2=\frac12$ and $V_3=0$ we have
$d_1=d_2=\sqrt{6c_1-2}$, and one can ascertain that we indeed
obtain real solutions for the specified parameter values, namely,
$\A_{1}=-\A_{2}=
-\frac{\sqrt{c_1}}{2}+\frac{\sqrt{c_3}}{2}\e{1}-\frac12\e{2} +
\sqrt{c_2}\e{3} +\frac12\e{13}+
\frac{\sqrt{c_3}}{2}\e{23}-\frac{\sqrt{c_1}}{2}\e{123}$ with
$c_3=5-\sqrt{5}$.

It should be noted, however, that the derived symbolical
expressions do not guarantee that we can always find {\it real}
parameters $V_2$ and $V_3$ which are required to ensure real
square root of MV. For example, if instead of the above MV we
would try to find square root of $\B=-\e{3}+\e{12}$, which was
used in the Example~2 of \cl{3}{0} algebra, we would find
$s_1=\tfrac{1}{2}$ and $s_2=-\tfrac{1}{2}$. The first value then
yields  $v_1=-V_1=\tfrac{1}{2} \sqrt{-4 V_2^2-(1+2 V_3)^2}$,
$v_2=-V_2$, $v_3=-1-V_3$ and for the second we obtain
$V_1=-v_1=\tfrac{1}{2} \sqrt{-4 V_2^2-(1-2 V_3)^2}$, $v_2=-V_2$,
and $v_3=1-V_3$. Taking the square of symbolical expressions one
can easily check that we indeed have the MV $-\e{3}+\e{12}$. It is
obvious, however, that in both cases the expression under square
root always is negative. Therefore, in fact, there are no real
square root of $-\e{3}+\e{12}$. We do not explore this question in
detail, since for explicitly given symbolic expressions the
analysis is straightforward although the answer is lengthy.



\section{\cl{1}{2} algebra}\label{sec:Cl12}

\subsection{The generic case, $s^2 \neq S^2$}

The analysis in this algebra is similar to \cl{3}{0}. The
coefficients $v_i$ and $V_i$ are
\begin{align}
v_1&=\frac{b_1s+b_{23}S}{2(s^2+S^2)},
   &v_2&=\frac{b_{2}s+b_{13}S}{2(s^2+S^2)},
    &v_3&=\frac{b_3 s -b_{12}S}{2(s^2+S^2)}, \label{vV1Cl12} \\
V_1&=\frac{b_{23}s-b_{1}S}{2(s^2+S^2)},
&V_2&=\frac{b_{13}S-b_{2}S}{2(s^2+S^2)},
&V_3&=-\frac{b_{12}s+b_{3}S}{2(s^2+S^2)}.\label{vV2Cl12}
\end{align}
Using the same replacement \eqref{sysASCl30} as  for \cl{3}{0}
case and introducing new  $b_{S}$ and $b_{I}$ (with
$D=b_{S}^2+b_{I}^2$),
\begin{equation}\label{bSbICl12}
\begin{split}
&b_{S}= \langle \B \cliffordconjugate{\B}\rangle_0 =
b_{0}^2-b_{1}^2+b_{2}^2+b_{3}^2-b_{12}^2-b_{13}^2+b_{23}^2-b_{123}^2,\\
&b_{I}= \langle \B \cliffordconjugate{\B} I \rangle_0 = 2 b_{3}
b_{12}-2 b_{2} b_{13}+2 b_{1} b_{23}-2 b_{0} b_{123}.
\end{split}
\end{equation}
we can write down the equations for scalar and pseudoscalar parts
in the form
\begin{equation}\label{sysARCl12}
\begin{split}
&2(4 T + b_{0}) (4 t -b_{123}) - b_{I}=0,\\
  &(b_0-b_{123}+4 T+4 t) (b_{0}+b_{123}+4 T-4 t)-b_{S}=0,
\end{split}\end{equation}
which reproduces the Eq.~\eqref{sysARCl30} for \cl{3}{0} case,
except that explicit form of $b_S$ and $b_I$ now is different. The
solutions for scalars $s$ and $S$ are given by \eqref{solACl30}
and \eqref{solBCl30}.  In the case $b_I=0$, the solution
\eqref{solBCl30} should be replaced by \eqref{solBd01Cl30}. After
insertion of  $s$ and $S$  into Eq.~\eqref{vV1Cl12},
\eqref{vV2Cl12} and  the result into \eqref{mvA30} yields four
real square roots of $\B$.


\subsection{The special case $s^2 = S^2$}

\subsubsection{The subcase $s=S\neq 0$}
When $s=S$ instead of \eqref{sysAspec1Cl30} we have very similar
compatibility conditions
\begin{equation}\begin{split}\label{sysAspec1Cl12}
  (4 s^2 b_0-b_3 b_{12}+b_2 b_{13}-b_1 b_{23})=0,\\
  (16 s^4-b_1^2+b_2^2+b_3^2-b_{12}^2-b_{13}^2+b_{23}^2-8 s^2 b_{123})=0,
 \end{split}\end{equation}
as in the case of \cl{3}{0}. When the shortcuts $b_I$ and $b_S$
are introduced the conditions again exactly coincide with
\eqref{sysAspec30}. Therefore, the solutions expressed in terms of
$b_I$ and $b_S$ are given by \eqref{solSspec1Cl30} and formulas
\eqref{vV1Cl12}, \eqref{vV2Cl12} and \eqref{mvA30}.

\subsubsection{The subcase $s=-S\neq 0$}
When $s$ is rewritten in terms of $b_I$ and $b_S$, we again can
use the same formulas as in the case of \cl{3}{0}.

\subsubsection{The subcase $s=S= 0$}
\cl{3}{0} formula \eqref{solSspec3Cl30} remains valid. Of course,
explicitly solved expressions for $v_1$ and $V_1$ now differ from
\cl{3}{0} case, because for \cl{1}{2} algebra we have  $v^2=v_1^2
- v_2^2 - v_3^2$ instead of $v^2=v_1^2 + v_2^2 + v_3^2$. In
particular,  the solution of \eqref{solSspec3Cl30} in this algebra
is
\begin{equation}
\begin{split}
&v_{1}=\pm \frac{c_1}{\sqrt{2}},\qquad
  V_{1}=\pm \frac{1}{c_1} \frac{b_{123}+2 (v_{2} V_{2}+v_{3} V_{3})}{\sqrt{2}},\qquad \textrm{with}\quad \\[3pt]
&c_1=\Bigl(\pm\sqrt{\left(b_{0}+v_{2}^2+v_{3}^2-V_{2}^2-V_{3}^2\right)^2+
   (b_{123}+2 (v_{2} V_{2}+v_{3} V_{3}))^2}\\
&\phantom{\frac{1}{\sqrt{2}}b_{0}-v_{2}^2-v_{3}^2  \big(b_{123}-2
(v_{2} V_{2}+v_{3} V_{3})\big)x}
 +b_{0}+v_{2}^2+v_{3}^2-V_{2}^2-V_{3}^2\Bigr)^{\frac{1}{2}},
\end{split}
\end{equation}
and can be compared to the similar expression
\eqref{exampleCl30free} for \cl{3}{0}. The remaining coefficients
$v_2,v_3,V_2,V_3$  may have arbitrary values and can be
treated as free parameters.

\subsection{Examples for $\cl{1}{2}$}
\subsubsection*{Example~7. The case $s^2\neq S^2$.}
Using the same initial MV  $\B=\e{1}-2 \e{23}$ as in
Example~1 we obtain the same values for $b_I,
b_S$ and $s, S, v_i, V_i$. After substituting them into
\eqref{mvA30} the square roots are found to be
\begin{equation}
\begin{split}
  &\sqrt{\B_{1,2}}=\A_{1,2}=\pm\frac12 (c_2 (-\e{1}+\e{1 2 3})-c_1 (1+\e{2 3})),\\
  &\sqrt{\B_{3,4}}=\A_{3,4}=\pm \frac12 (-c_1 (\e{1}+\e{1 2 3})+c_2 (-1+\e{2 3})),
\end{split}\end{equation}
where $c_1=\sqrt{-2+\sqrt{5}}$ and  $c_2=\sqrt{2+\sqrt{5}}$\,.


\section{\cl{2}{1} algebra}\label{sec:Cl21} In~\cite{Marchuk2015}, considering the
CA representation as a tensor product it is shown that ``for odd
$n \ge 3$, there are three classes of isomorphic Clifford
algebras... [what] is consistent with Cartan's classification of
real Clifford algebras.'' Namely, \cl{3}{0} and \cl{1}{2}  are
represented by $2\times 2$ complex matrices $\bbC(2)$. The
similitude of expressions  for square roots in \cl{3}{0} and
\cl{1}{2} algebras confirms that these two algebras fall into the
same class. Algebras \cl{0}{3} and \cl{2}{1} are represented by
blocked $2\times 2$ and $1\times 1$ matrices ${^2}\bbR(2)$ and
${^2}\bbH(1)$, respectively.  Therefore, they belong to different
classes. Indeed, as we shall show below, the analysis for
\cl{2}{1} is more similar to \cl{0}{3}. However, in this case
there appear significant and interesting differences since
\cl{0}{3} and \cl{2}{1} belong to different classes.


\subsection{The generic case, $s^2 \neq S^2$}
The system of nonlinear equations is
\begin{align}
b_0&=s^2+S^2+\bv^2+\bV^2, &b_{123}&=2(sS+\bv\cdot\bV)\label{MVM10Cl21}, \\
b_1&=2(s v_1+S V_1), &b_{23}&=2(s V_1+S v_1)\label{MVM20Cl21}, \\
b_2&=2(s v_2+S V_2), &b_{13}&=-2(s V_2+S v_2)\label{MVM30Cl21}, \\
b_3&=2(s v_3+S V_3), &b_{12}&=-2(s V_3+S v_3)\label{MVM40Cl21},
\end{align}
where now $\bv^2=v_1^2+v_2^2-v_3^2$ and $\bv\cdot\bV=v_1 V_1 + v_2 V_2- v_3 V_3$. When $s^2-S^2\neq
0$ the solution of Eqs.~\eqref{MVM20Cl21}-\eqref{MVM40Cl21} is
\begin{align}
v_1&=\frac{b_1s-b_{23}S}{2(s^2-S^2)}, & v_2&=\frac{b_{2}s+ b_{13}S}{2(s^2-S^2)},&  v_3&=\frac{b_3s+b_{12}S}{2(s^2-S^2)},\label{v123Cl21}\\
V_1&=\frac{b_{23}s-b_{1}S}{2(s^2-S^2)},
&V_2&=-\frac{b_{13}s+b_{2}S}{2(s^2-S^2)},
&V_3&=-\frac{b_{12}s+b_{3}S}{2(s^2-S^2)}.\label{V123Cl21}
\end{align}
Insertion  of $v_i$ and $V_i$ into \eqref{MVM10Cl21} gives two
coupled equations for unknowns $s, S$
\begin{equation}\label{sysA0Cl21}
\begin{split}
& b_{S}+4 s^2 (-6 S^2+b_0)+8 s S b_{123}=4 s^4+(-2 S^2+b_0)^2+b_{123}^2,\\
  & b_{I}=2 (2 (s^2+S^2)-b_0) (4 s S-b_{123}),
\end{split}\end{equation}
  where $b_{S}$ and $b_{I}$ are functions of coefficients of MV $\B$,
\begin{equation}\label{bSbICl03A}
\begin{split}
&b_{S}= \langle \B \cliffordconjugate{\B}\rangle_0 =
b_{0}^2-b_{1}^2-b_{2}^2+b_{3}^2+ b_{12}^2-b_{13}^2-b_{23}^2+b_{123}^2,\\
&b_{I}= \langle \B
\cliffordconjugate{\B} I \rangle_0 = -2 b_{3} b_{12}+2 b_{2} b_{13}-2 b_{1} b_{23}+2 b_{0} b_{123}.
\end{split}
\end{equation}
Since the Eqs.~\eqref{sysA0Cl21} and \eqref{V123Cl03} have
  the same shape (the definitions of $b_{S}$ and $b_{I}$ in these
  equations, of course, differ) we can make use of \eqref{sysAS0Cl03} to
  lower the order of the system. However, there comes one
important difference: in \cl{2}{1} the determinant norm,
$D=b_S^2-b_I^2$, is not always positive. It can happen that for a
some $\B$ the norm becomes negative, $D<0$. In this case the
solution set must be considered empty. The other particularity is
that in the solution \eqref{solB0Cl03} instead of single
($-\sqrt{D}$) we have to take into account both signs, i.e.,
$\pm\sqrt{D}$, what doubles the number of possible solutions in the
case $D>0$,
\begin{equation}\label{solB0Cl21}
\begin{split}
&t_{1,2,3,4}=\frac{1}{4} \Bigl(b_{123} \pm
\frac{1}{\sqrt{2}}\sqrt{b_{S}\pm\sqrt{D}}\Bigr),\quad
T_{1,2,3,4}=\frac{1}{4}\Bigl(\frac{\pm
b_{I}}{\sqrt{2}\sqrt{b_{S}\pm\sqrt{D}}}+ b_{0}\Bigr).
\end{split}
\end{equation}
Here again we can take signs for $t$ in all possible combinations,
but then the signs for $T$ must follow the same upper-lower
positions as in  $t$. When  $b_{I}=0$, or equivalently
$b_{S}\pm\sqrt{D} = 0$, the solution set is
\begin{equation}\label{solB0d01aCl21}
\begin{split}
(t_{1,2},\ T_{1,2})= \bigl(\tfrac{1}{4}b_{123},\ \tfrac{1}{4}
(\pm\sqrt{b_{S}}+b_{0})\bigr).
\end{split}
\end{equation}
Since we also have four sign combinations in the solution for
$s,S$ (as in \eqref{solA0Cl03}), we end up with 16~different
square roots of MV in a generic case of \cl{2}{1}.

\subsection{The special case $s^2 = S^2$}

The analysis again closely follows \cl{0}{3} case, except that we
have different signs in the expressions.
\subsubsection{The subcase $s=S\neq 0$}
Now the coefficients satisfy the conditions $b_1 =b_{23}$, $b_2 =
-b_{13}$,  $b_3= - b_{12}$.  This  allows to eliminate the
singularity at $s=S$, and as a result the system of
Eqs.~\eqref{MVM20Cl21}-\eqref{MVM40Cl21} has a special solution,
\begin{equation}\begin{split}\label{vVSpecialCl21}
v_{1}= \frac{b_{1}}{2 s}-V_{1},\quad v_{2}= \frac{b_{2}}{2 s}-V_{2},\quad v_{3}= \frac{b_{3}}{2 s}-V_{3},
\end{split}\end{equation}
which exactly coincides with the same solution for \cl{0}{3} (see
Eq.~\eqref{vVSpecialCl03}). Thus, after doing similar computations
we find that Eq.~\eqref{AspecCl03} is to be replaced by
\begin{equation}\begin{split}\label{AspecCl21}
  V_{1}=& \frac{\sqrt{2}}{8 s}\Bigl(\sqrt{2} b_{1}\pm \bigl(8 s^2 \left(b_{0}-b_{123}+4 (-V_{2}^2+V_{3}^2)\right)+16 s (b_{2} V_{2}-b_{3} V_{3})\\
  &\phantom{\bigl(-8 s^2 \left(b_{0}-b_{123}+4 (V_{2}^2+V_{3}^2)\right)}
  -b_{0}^2+2 b_{1}^2-b_{123}^2+b_{S}\bigr)^{1/2}\Bigr).
\end{split}\end{equation}
The coefficients $s_{1}$ and   $s_{2}$ are very similar to
\eqref{solSspec1Cl03}, except that now we need to take into
account all passible sign combinations,
\begin{equation}\begin{split}\label{solSspec1Cl21}
s_{1,2} =& \pm\tfrac{1}{2\sqrt{2}}\sqrt{\pm\sqrt{2
b_{S}-(b_{0}-b_{123})^2}+b_{0}+b_{123}}.
\end{split}\end{equation}
The listed formulas  solve the square root problem in the case
$s=S\neq 0$.

\subsubsection{The subcase $s=-S\neq 0$}
The only formulas which differ from those in \cl{0}{3} algebra are
related with the coefficient compatibility condition  $b_1 =
-b_{23}$, $b_2 = b_{13}$, $ b_3= b_{12}$. Now, they should
be replaced by
\begin{equation}\begin{split}\label{Aspec2Cl21}
  V_{1}=& \frac{\sqrt{2}}{8 s}\Bigl(-\sqrt{2} b_{1}\pm \bigl(8 s^2 \left(b_{0}+b_{123}+4(-V_{2}^2+V_{3}^2)\right)\\
  &\phantom{\bigl(-8 s^2}-16 s (b_{2} V_{2}-b_{3} V_{3})-b_{0}^2+2 b_{1}^2-b_{123}^2+b_{S}\bigr)^{1/2}\Bigr),
\end{split}\end{equation}
%
\begin{equation}\begin{split}\label{solSspec1NegCl21}
s_{1,2} =& \pm\tfrac{1}{2\sqrt{2}}\sqrt{\pm\sqrt{2
b_{S}-(b_{0}+b_{123})^2}+b_{0}-b_{123}}\,.
\end{split}\end{equation}
The remaining formulas that are required to write down the
answer exactly match the formulas in the corresponding subsection
of \cl{0}{3} algebra.


\subsubsection{The subcase $s=S= 0$}
The only distinct formulas from \cl{0}{3} case are listed below,
\begin{equation}\begin{split}
&v_{1}=\pm \frac{c_1}{\sqrt{2}},\qquad
  V_{1}=\pm \frac{1}{c_1} \frac{b_{123}+2 (-v_{2} V_{2}+v_{3} V_{3})}{\sqrt{2}},\qquad \textrm{where}\quad \\[3pt]
&c_1=\Bigl(\pm\sqrt{\left(b_{0}-v_{2}^2+v_{3}^2-V_{2}^2+V_{3}^2\right)^2-
   (b_{123}+2 (-v_{2} V_{2}+v_{3} V_{3}))^2}\\
&\phantom{\frac{1}{\sqrt{2}}b_{0}-v_{2}^2-v_{3}^2  (b_{123}-2
(v_{2} V_{2}+v_{3} V_{3}))}
+b_{0}-v_{2}^2+v_{3}^2-V_{2}^2+V_{3}^2\Bigr)^{\frac{1}{2}} .
\end{split}\end{equation}
This ends the investigation of the square root formulas for all
real Clifford algebras in 3D.

\subsection{Examples for $\cl{2}{1}$}
\subsubsection*{Example~8. The case $s^2\neq S^2$.}
First we shall show  that MV  $\B=\e{1}-2 \e{23}$  has no real
square roots. Indeed, we have  $b_S=-5$,  $b_I=4$ and
$D=b_S^2-b_I^2=(3)^2$. As a result  the expression under square
root in \eqref{solB0Cl21}, namely $b_S\pm\sqrt{D}=-5\pm3$, is
always negative and there are no real valued solutions.

Next we shall calculate the roots of $\B=2+\e{1}+\e{13}$. The
values of $b_I, b_S$ are $0$ and $2$, respectively. The
determinant norm of the MV is positive, $D=4>0$. Using
\eqref{solB0Cl21} we find four real values for pairs $(T_i,t_i)$.
In particular we obtain
 $ (T_1,t_1)=\bigl(\frac{1}{4} (2-\sqrt{2}) , 0\bigr)$, $(T_2,t_2)=\bigl(\frac{1}{4} (2+\sqrt{2}) , 0\bigr)$, $(T_3,t_3)=\bigl(\frac{1}{2} , -\frac{1}{2 \sqrt{2}}\bigr)$ and $(T_4,t_4)=  \bigl(\frac{1}{2} , \frac{1}{2 \sqrt{2}}\bigr)$.
Inserting them into \eqref{solA0Cl03} we obtain 16 pairs
$(s_i,S_i)$ of scalars:
\[ \begin{array}{lll}
 (s_1= 0,S_1=-\frac{c_2}{\sqrt{2}}),
  &\mspace{-15mu}(s_2= 0,S_2=\frac{c_2}{\sqrt{2}}),
  &\mspace{-15mu}(s_3= -\frac{c_2}{\sqrt{2}},S_3= 0),\\
  (s_4= \frac{c_2}{\sqrt{2}},S_4= 0),
  &\mspace{-15mu}(s_5= 0,S_5= -\frac{c_1}{\sqrt{2}}),
  &\mspace{-15mu}(s_6= 0,S_6= \frac{c_1}{\sqrt{2}}),\\
  (s_7= -\frac{c_1}{\sqrt{2}},S_7= 0),
  &\mspace{-15mu}(s_8= \frac{c_1}{\sqrt{2}},S_8= 0),
  &\mspace{-15mu}(s_9= -\frac{c_2}{2},S_9= \frac{c_1}{2}),\\
  (s_{10}= \frac{c_2}{2},S_{10}= -\frac{c_1}{2}),
  &\mspace{-15mu}(s_{11}= -\frac{c_1}{2},S_{11}= \frac{1}{\sqrt{2} c_1}),
  &\mspace{-15mu}(s_{12}= \frac{c_1}{2},S_{12}= -\frac{1}{\sqrt{2} c_1}),\\
  (s_{13}= -\frac{c_2}{2},S_{13}= -\frac{c_1}{2}),
  &\mspace{-15mu}(s_{14}= \frac{c_2}{2},S_{14}= \frac{c_1}{2}),
  &\mspace{-15mu}(s_{15}= -\frac{c_1}{2},S_{15}= -\frac{1}{\sqrt{2} c_1}),\\
  (s_{16}= \frac{c_1}{2},S_{16}=\frac{1}{\sqrt{2} c_1}), & &
  \end{array}\]
where $c_1=\sqrt{2+\sqrt{2}}$ and $c_2=\sqrt{2-\sqrt{2}}$. After
substitution of pairs $(s_i,S_i)$ into Eq.~\eqref{v123Cl21} and
then into Eq.~\eqref{mvA30} we obtain 16  roots
$\sqrt{\B_{i,j}}=\A_{i,j}$ the squares of which give the initial
MV $\B=2+\e{1}+\e{13}$,
\begin{equation}
\begin{split}
  &\A_{1,2}=\pm\tfrac{1}{2}\big(c_1 \e{2}-c_1 \e{2 3}-\sqrt{2}c_2\e{123}\big),\\
 &\A_{3,4}=\pm\tfrac{1}{\sqrt{2}}\big(-c_1^{-1}\e{2}+c_1^{-1}\e{2 3}+c_1 \e{1 2 3}\big),\\
&\A_{5,6}=\pm\tfrac{1}{2}\big(\sqrt{2}c_2+c_1\e{1}+c_1\e{13}\big),\\
&\A_{7,8}=\pm\tfrac{1}{\sqrt{2}}\big(c_1+ c_1^{-1}\e{1}+c_1^{-1}\e{1 3}\big),\\
&\A_{9,10}=\pm\tfrac{1}{2\sqrt{2}}\big(\sqrt{2}c_2-c_2\e{1}-c_1\e{2}-c_2\e{13}+c_1\e{2 3}+\sqrt{2}c_1 \e{1 2 3}\big),\\
  &A_{11,12}=\pm\tfrac{1}{2\sqrt{2}}\big(\sqrt{2}c_ 1+c_1 \e{1}-c_ 2 \e{2}+c_1 \e{1 3}+c_2 \e{2 3}-2c_1^{-1}\e{1 2 3}\big),\\
&\A_{13,14}=\pm\tfrac{1}{2\sqrt{2}}\big(\sqrt{2}c_1+c_1 \e{1}+c_2 \e{2}+c_1 \e{1 3}-c_2 \e{2 3}+2c_1^{-1}\e{1 2 3}\big),\\
&\A_{15,16}=\pm\tfrac{1}{2\sqrt{2}}\big(-\sqrt{2}c_2+c_2 \e{1}-c_1
\e{2}+c_2 \e{13}+c_1 \e{2 3}+\sqrt{2}c_1 \e{1 2 3}\big) .
\end{split}\end{equation}


\section{Application: Riccati equation in Clifford algebra\label{sec:application}}
Solution of quaternionic Riccati equation, which is quadratic with
respect to unknown variable, has been presented
in~\cite{Janovska13}. The simplest quadratic equations of this
type in CA for unknown MV $\X$ are
\begin{equation}\label{RiccA}
\X^2=\B\quad\text{and}\quad\X\A\X=\B,
\end{equation}
where $\A$ and $\B$ are constant MVs. The solution of
Eq.~\eqref{RiccA} are expressed through the square root and
inversion of MV,
\begin{equation}\label{RiccB}
\X=\pm\sqrt{\B}\quad\text{and}\quad\X=\pm\sqrt{\B\A}\;\A^{-1}=\pm\A^{-1}\sqrt{\A\B}\,.
\end{equation}
More general Riccati-Clifford equation contains linear term,
\begin{equation}\label{RiccC}
\X\A\X+\C\X+\X\C=\B.
\end{equation}
The equation can be easily solved if $\C$ belongs to CA center,
i.e., $\C=\alpha+\beta I$, where $\alpha,\beta\in\bbR$. The
solution is~\cite{Dargys2019}
\begin{equation}\label{RiccD}
\X_{1,2}=(-\C\pm\sqrt{\B\A+\C^2}\,)\A^{-1}.
\end{equation}
If MVs are replaced by scalars ($\A\to a$, $\B\to c$, $\C\to b/2$)
and the geometric product is treated as standard scalar product
then Eq.~\eqref{RiccD} reduces to the well-known scalar quadratic
equation for roots $x_{1,2}=(-b\pm\sqrt{b^2+4ac})/2a$. The
principle difference between solutions of Riccati-Clifford and
standard quadratic equation lies in a number of roots. In the latter
case there may be at maximum two different roots whereas in the case
of MV the number of roots depends on algebra and structure of the
MV itself.

Our work is a step forward in a solution of  more general
quadratic equations in Clifford algebras. Indeed, if a quadratic
equation for a MV can be factorized into quite a general form $(\A
\X +\X\B +\C)^2=\D$ then using the presented algorithm to find
$\sqrt{\D}$ and the explicit solution of Sylvester equation in
CA~\cite{Dargys2018} for $n\le 3$ one can immediately find all
solutions of the quadratic MV equation.


\section{Conclusions\label{sec:conclusions}}
First, we have  shown that the  square root of general MV in
\cl{3}{0}, \cl{0}{3}, \cl{1}{2} and \cl{2}{1} algebras can be
expressed in radicals.  A detailed algorithm for all four algebras
is described which can be used to find the roots of MVs in a
closed form. In our first paper on the subject~\cite{Dargys2019},
we considered the roots of individual grades and provided explicit
formulas in a coordinate free form. However, for a general MV the
latter form, where many condition are controlled by plus/minus
signs, appears rather complicated.

Second, it is shown that apart from two standard (plus/minus)
  roots there may appear more (up to 16 for \cl{2}{1}) isolated roots in CA, or
the roots may not exist at all. The MV roots in the CA may
accommodate a number free parameters that bring a continuum of
square roots on respective parameter hypersurfaces and therefore
become an important ingredient in the nonlinear CA problems.

Third, the Gr{\"o}bner basis algorithm was applied in the
algorithm analysis using \textit{Mathematica}
commands~\cite{AcusDargys2017}. We also implemented purely
numerical search of square roots using {\it Mathematica} universal
root search algorithm realized in the system function
\textbf{FindInstance[~]}. This was done in hope to find some
isolated roots which our algorithm might miss. We didn't find any
new numerical roots. On the contrary, the command
\textbf{FindInstance[~]} often -- especially in the case $s^2
=S^2\neq 0$ -- was unable to find any roots at all, whereas our
algorithm worked flawlessly. The only complication which we
encountered in the algorithm implementation was that {\it
Mathematica} symbolic zero detection algorithm
\textbf{PossibleZeroQ[~]} in more complicated cases often had to
switch to numerical procedure to conclude that complicated
symbolic expression indeed represents zero. This is quite
understandable, since it is well known that two expression
equivalence problem is, in general, undecidable.

Fourth, we found that for algebras $\cl{3}{0}$ and $\cl{1}{2}$ the
square root solution in general is a union of following sets: 1)
when $s^2\neq S^2$, we have a set of (up to) four isolated roots,
2) when $s^2 =S^2\neq 0$, we have a set of two roots  and 3) in
addition, we may have a set of infinite number of roots that
belong to at most a four dimensional parameter manifold. For case
of $\cl{0}{3}$ algebra the square root is a union of  sets: 1) a
set of (up to) four isolated roots in the case $s^2\neq S^2$, 2)
an infinite number of roots that make at most two dimensional
manifold in the case $s^2 =S^2\neq 0$, and 3) a set of infinite
number of roots which make a at most four dimensional manifold.
For $\cl{2}{1}$ algebra, the square root is a union of following
sets: 1) in  the case $s^2\neq S^2$, a set of (up to) sixteen
isolated roots, 2) in the case $s^2 =S^2\neq 0$, a infinite number
of roots which belong at most to two dimensional manifold and 3) a
set of infinite number of roots that belongs at most to 4D
parameter manifold. Since the solutions of each of different
(sub)sets have different values of scalars $s$ and $S$ in the
parameter spaces that follow from Eq.~\eqref{mvA30}, the solution
in one subset cannot be obtained from solution in a different
subset (even from the infinite one). The algorithm should also
work for the case of complex Clifford algebras with the only
simplification that in this case there is no need to look
for a positivity of expressions inside the square roots.

The proposed algorithm is a step forward in a solution of general
quadratic equations in Clifford algebras and may find application
in control and systems theory~\cite{Abou2003}, since the CA
solutions reflect totally new properties of MV roots, for example,
the root multiplicity and appearance of free parameters.  Due to
complexity of root algorithms in CA, it is preferred to do such
calculations with in advance prepared numerical/symbolic
subroutines. Nowadays,  nobody is worried how to calculate the
square root of a number with a pocket calculator in hands.
Probably, with a fast development of CA the history will repeat
itself once more but now with a multivector square root.



% ----------------------------------------------------------------
%\paragraph{Acknowledgment.}
% ----------------------------------------------------------------
%We would like to thank ...

%\begin{thebibliography}{10}


%\end{thebibliography}


%\bibliographystyle{REPORT}
%\bibliography{SqrtRootMV}
\begin{thebibliography}{10}

\bibitem{Grant2015}
H.~Grant and I.~Kleiner,
\newblock {\em Turning Points in the History of Mathematics},
\newblock Springer, New York, 2015.

\bibitem{Cayley1872}
A.~Cayley,
\newblock On the extraction of square root of matrix of the third order,
\newblock Proc. Roy. Soc. Edinburgh {\bf 7}, 675--682 (1872).

\bibitem{Higham08}
N.~J. Higham,
\newblock {\em Functions of Matrices (Theory and Computation)},
\newblock SIAM, Philadelphia, 2008.

\bibitem{Niven1942}
I.~Niven,
\newblock The roots of quaternion,
\newblock The American Mathematical Monthly {\bf 49}(6), 386--388 (1942).

\bibitem{Opfer2017}
G.~Opfer,
\newblock {N}iven's algorithm applied to the roots of the companion polynomial
  over $\mathbb{R}^4$ algebras,
\newblock Adv. Appl. Clifford Algebras {\bf 27}, 2659--2675 (2017).

\bibitem{Falcao2018}
M.~I. Falc{\~a}o, F.~Miranda, R.~Severino, and M.~J. Soares,
\newblock On the roots of coquaternions,
\newblock Adv. Appl. Clifford Algebras {\bf 28}, 97 (2018).

\bibitem{Ozdemir2009}
M.~{\"O}zdemir,
\newblock The roots of a split quaternion,
\newblock Appl. Math. Lett. {\bf 22}, 258--263 (2009).

\bibitem{Sangwine2006}
S.~J. Sangwine,
\newblock Biquaternion (complex quaternion) roots of $-1$,
\newblock Adv. Appl. Clifford Algebras {\bf 16}(1), 63--68 (2006).

\bibitem{Chapell2015}
J.~M. Chappell, A.~Iqbal, L.~J. Gunn, and D.~Abbott,
\newblock Functions of multivector variables,
\newblock PLoS ONE {\bf 10}(3), 1--21 (2015),
\newblock Doi:10.1371/journal.pone.0116943.

\bibitem{Hitzer2011}
E.~Hitzer and R.~Ab{\l}amowicz,
\newblock Geometric roots of $-1$ in {C}lifford algebras \textit{Cl}$_{p,q}$
  with $p+q\le 4$,
\newblock Adv. Appl. Clifford Algebras {\bf 21}(1), 121--144 (2011).

\bibitem{Hitzer2013}
E.~Hitzer, J.~Helmstetter, and R.~Ab{\l}amowicz,
\newblock Square roots of $-1$ in real {C}lifford algebras,
\newblock in {\em Quaternion and Clifford-Fourier Transforms and Wavelets},
  edited by E.~Hitzer and S.~J. Sangwine, pages 123--153, Springer, Basel,
  2013.

\bibitem{Hitzer13B}
E.~Hitzer and {S.~J. Sangwine, (Eds.)},
\newblock {\em Quaternion and Clifford-Fourier Transforms and Wavelets},
\newblock Springer, Basel, 2013.

\bibitem{Dargys2019}
A.~Dargys and A.~Acus,
\newblock Square root of a multivector in 3D {C}lifford algebras,
\newblock Nonlinear Analysis: Modelling and Control , 1--20 (2019),
\newblock In press.

\bibitem{AcusDargys2017}
A.~Acus and A.~Dargys,
\newblock Geometric Algebra Mathematica package, 2017,
\newblock \url{https://github.com/ArturasAcus/GeometricAlgebra} , 65,
\newblock The full version is to be described elsewhere.

\bibitem{Doran03}
C.~Doran and A.~Lasenby,
\newblock {\em Geometric Algebra for Physicists},
\newblock Cambridge University Press, Cambridge, 2003.

\bibitem{Lounesto97}
P.~Lounesto,
\newblock {\em Clifford Algebra and Spinors},
\newblock Cambridge University Press, Cambridge, 1997.

\bibitem{Acus2018}
A.~Acus and A.~Dargys,
\newblock The inverse of multivector: Beyond the threshold $p+q=5$,
\newblock Adv. Appl. Clifford Algebras {\bf 28}, 65 (2018).

\bibitem{Marchuk2015}
N.~G. Marchuk,
\newblock Demonstration representation and tensor products of {C}lifford
  algebras,
\newblock {T}rudy {M}atematicheskogo {I}nstituta {I}meni {V.A.} {S}teklova {\bf
  290}, 154--165 (2015),
\newblock English translation: Proc. Steklov Inst. Math., Pleiades Publishing,
  Ltd., {\bf 290}, 143-154 (2015).

\bibitem{Janovska13}
D.~Janovsk{\'a} and G.~Opfer,
\newblock The algebraic {R}iccati equations for quaternions,
\newblock Adv. Appl. Clifford Algebras {\bf 23}, 907--918 (2013).

\bibitem{Dargys2018}
A.~Dargys and A.~Acus,
\newblock A note on solution of ax+xb=c by Clifford algebras,
\newblock arXiv (1902.09194 [math.RA]), 1--4 (2018).

\bibitem{Abou2003}
H.~Abou-Kandil, G.~Freiling, V.~Ionescu, and G.~Jank,
\newblock {\em Matrix Riccati Equations (in Control and Systems Theory)},
\newblock Birkh{\"a}user Verlag, Basel, 2003.

\end{thebibliography}




%\bibitem{Dargys2018}
%A. Dargys and A. Acus,
%\newblock A note on solution of $ax+xb=c$ by Clifford algebras, 2017,
%\newblock arXiv:1902.09194  [math.RA] , 1--4 (2018).

\end{document}

